\documentclass[10pt,a4paper]{article}
\usepackage[margin=2cm,headheight=26pt,headsep=10pt]{geometry}
\usepackage{times}
\usepackage[hidelinks]{hyperref}
\usepackage{enumitem}
\setlist{nosep,leftmargin=1.2em}
% Page header required by the course on every page
\usepackage{fancyhdr}
\newcommand{\groupno}{XX}
\newcommand{\projname}{INN-RC: Lossless Image Steganography with Message-Conditioned Integer Coupling}
\fancypagestyle{uct}{%
  \fancyhf{}%
  \fancyhead[L]{\footnotesize UCT702 Research Project \quad Group No.: \groupno \\ Project Name: \projname}%
  \fancyhead[R]{\footnotesize Reetam Dutta (102318024)\\ Priyadarshi Vivekanand (102318048)}%
  \fancyfoot[C]{\footnotesize \thepage}%
  \renewcommand{\headrulewidth}{0.4pt}}

\pagestyle{uct}
\setlength{\parskip}{3pt}
\begin{document}
\thispagestyle{uct}
\begin{center}
{\large\bfseries Phase 1: Literature Review and Problem Statement}\\[2pt]
Track TRK-07: Lossless End-to-End Steganography (Deep Learning, Computer Vision)
\end{center}

\section*{1. Why this track}
We picked TRK-07 because it is the only track whose toolbox column lists Computer Vision together with Deep Learning,
which is the stack we wanted to work with. The task is to hide a message in an image with a neural network and still
give back both the message and the cover image without a single wrong bit. That last condition is what makes the
problem hard: neural networks work with real numbers, images are stored as 8-bit integers, and the rounding in
between destroys exactly the information that reversibility needs.

\section*{2. Baseline paper}
The project PDF names ``Deep Reversible Image Watermarking Network via Invertible Residual Compensation'' for TRK-07,
but the link given for it opens a different article (Saeidi and Mirzakuchaki, \emph{Reversible data hiding in
encrypted images using LWE-based secret sharing}, Scientific Reports, 2025), which uses no neural network. We could not
find a paper with the listed title. The closest published work that matches the description (invertible network,
auxiliary data, exact recovery) is \textbf{Deep Robust Reversible Watermarking (DRRW)} by Chen \emph{et al.}
(arXiv:2503.02490, 2025) and its ICCV 2025 version \emph{CRMark}. We use DRRW as our baseline reference and noted the
link problem for the instructor.

\textbf{What DRRW does.} An integer invertible watermark network (iIWN) maps the pair (cover, watermark) to a stego
image and a latent variable. Because the network must stay invertible, the stego can leave the 0--255 range (overflow)
and the latent cannot be thrown away. Both are compressed with arithmetic coding into an auxiliary bitstream, which is
then hidden in the clipped stego with classical reversible data hiding (RDH). On a lossless channel the receiver
undoes the RDH layer, rebuilds the overflowed stego and the latent, and runs the network backwards. On a lossy channel
it simply reads the watermark, which survives noise and JPEG because of a noise layer used in training. The paper
adds an overflow penalty loss and reports a much shorter auxiliary bitstream than earlier robust-reversible schemes.

\textbf{Limitation we want to address.} In DRRW the latent variable always has to be stored, so part of the auxiliary
stream is spent on information that has nothing to do with the message. The networks are also large and need a GPU.

\section*{3. Related work}
\textbf{Classical RDH.} Difference expansion (Tian, 2003), histogram shifting (Ni \emph{et al.}, 2006) and
prediction-error expansion, PEE (Thodi and Rodr\'iguez, 2007), are the foundations. Sachnev \emph{et al.} (2009) added
the rhombus predictor, two-pass checkerboard embedding and sorting by local variance, and Ou \emph{et al.} (2013)
introduced pairwise PEE. These methods are exactly reversible but give no robustness: one changed pixel breaks the
extraction.

\textbf{CNN predictors for RDH.} Hu and Xiang replaced hand-made predictors with a CNN (IEEE SPL, 2021) and later
added adaptive embedding (IEEE TPAMI, 2022). Qiu \emph{et al.} (Entropy, 2025) also predict pixel complexity with the
network. Learning only improves the predictor; the embedding itself stays classical.

\textbf{Deep and invertible steganography.} Baluja (NeurIPS 2017) and HiDDeN (ECCV 2018) showed that networks can hide
data, HiDDeN with a noise layer for robustness. HiNet (ICCV 2021) and ISN (CVPR 2021) used invertible networks with
affine coupling (RealNVP style) so the same weights hide and reveal. RIIS (CVPR 2022), PRIS (2023) and Lan \emph{et al.}
(AAAI 2023) make the revealed data survive distortions. None of these return the cover exactly.

\begin{center}
\small
\begin{tabular}{lcccc}
\hline
Method & Learned & Exact cover recovery & Robust message & Extra data stored \\
\hline
PEE / rhombus PEE (Sachnev 2009) & no & yes & no & location map \\
CNN predictor RDH (Hu \& Xiang 2021/22) & predictor only & yes & no & location map \\
HiDDeN (2018) & yes & no & yes & -- \\
HiNet / ISN (2021) & yes (INN) & no & no & -- \\
RIIS / PRIS / Lan \emph{et al.} & yes (INN) & no & yes & -- \\
DRRW / CRMark (2025) & yes (integer INN) & yes & yes & latent + overflow \\
Ours (planned) & yes (integer coupling) & yes & yes & decoder errors + overflow \\
\hline
\end{tabular}
\end{center}

\section*{4. Problem statement}
Design a steganographic scheme, trained end-to-end, that (i) hides a binary message in an 8-bit grayscale image,
(ii) returns the message with BER $= 0$ and the cover bit-for-bit on a lossless channel, (iii) keeps the message
readable when the stego is mildly distorted, and (iv) keeps the auxiliary data needed for exact recovery small, so the
second (classical RDH) stage costs little image quality. We also want the model small enough to train on a CPU.

\textbf{Initial idea for the extension.} Instead of an invertible map that produces a latent, we plan to condition an
\emph{integer} additive coupling network on the message. Given the message, the cover is then recovered exactly by
subtraction, so no latent has to be stored. The only auxiliary data left are the bits the decoder reads wrongly and
the overflow pixels, which we code with an arithmetic coder whose contexts come from the decoder's own confidence.

\section*{5. Datasets and evaluation plan}
BOSSbase 1.01 could not be downloaded from our network, so we use BSD100 (100 images) and Kodak (24 images) as natural
image test sets, real DICOM CT/MR/CR scans from the public pydicom test collection as the medical set, and Urban100 for
training (disjoint from all test sets). Metrics: PSNR, SSIM, payload in bpp, BER, NC, exact-recovery rate, and BER under
Gaussian noise, salt-and-pepper noise and JPEG. NPCR/UACI are listed for the encrypted-domain tracks (TRK-03/04) and
do not apply to TRK-07.

\section*{6. Toolbox setup}
\begin{itemize}
\item MATLAB R2022b or newer, Image Processing Toolbox (image I/O, JPEG attack) and the App Designer components
(\texttt{uifigure}, \texttt{uiaxes}) for the GUI. Network inference is written with plain matrix code
(\texttt{conv3x3.m}), so the Deep Learning Toolbox is not needed to run the system.
\item The networks are trained in Python (PyTorch, CPU) and the weights are exported to \texttt{.mat} files that
MATLAB reads directly. Training scripts are in \texttt{training/}.
\end{itemize}

\section*{Key references}
\small
Chen \emph{et al.}, ``Deep robust reversible watermarking,'' arXiv:2503.02490, 2025.
Chen \emph{et al.}, ``Learning robust image watermarking with lossless cover recovery,'' ICCV 2025.
Sachnev \emph{et al.}, IEEE TCSVT 19(7), 2009. Hu and Xiang, IEEE SPL 28, 2021; IEEE TPAMI 44(12), 2022.
Jing \emph{et al.}, ``HiNet,'' ICCV 2021. Lu \emph{et al.}, ``ISN,'' CVPR 2021. Xu \emph{et al.}, ``RIIS,'' CVPR 2022.
Hoogeboom \emph{et al.}, ``Integer discrete flows,'' NeurIPS 2019. Zhu \emph{et al.}, ``HiDDeN,'' ECCV 2018.
\end{document}
